% =============================================================================
%  Bubble Radius Fitting — User Guide
%  Author: Zixiang (Zach) Tong, The University of Texas at Austin
% =============================================================================
\documentclass[11pt, letterpaper]{article}

% ---------- Packages ----------
\usepackage[margin=1in]{geometry}
\usepackage{graphicx}
\usepackage{booktabs}
\usepackage{enumitem}
\usepackage{xcolor}
\usepackage{hyperref}
\usepackage{fancyhdr}
\usepackage{titlesec}
\usepackage{amsmath}
\usepackage{listings}
\usepackage{float}
\usepackage{caption}
\usepackage{subcaption}
\usepackage{tcolorbox}
\usepackage{fontawesome5}

% ---------- Colour definitions ----------
\definecolor{primaryblue}{HTML}{3B82F6}
\definecolor{darknavy}{HTML}{0F172A}
\definecolor{successgreen}{HTML}{22C55E}
\definecolor{errorred}{HTML}{EF4444}
\definecolor{warningamber}{HTML}{FCD34D}
\definecolor{lightslate}{HTML}{F8FAFC}
\definecolor{bordergrey}{HTML}{E2E8F0}
\definecolor{dimgrey}{HTML}{94A3B8}

% ---------- Hyperlink style ----------
\hypersetup{
    colorlinks=true,
    linkcolor=primaryblue,
    urlcolor=primaryblue,
    citecolor=primaryblue,
}

% ---------- Header / footer ----------
\pagestyle{fancy}
\fancyhf{}
\fancyhead[L]{\small\textcolor{dimgrey}{Bubble Radius Fitting --- User Guide}}
\fancyhead[R]{\small\textcolor{dimgrey}{\thepage}}
\renewcommand{\headrulewidth}{0.4pt}

% ---------- Section styling ----------
\titleformat{\section}
  {\Large\bfseries\color{darknavy}}{\thesection}{1em}{}
\titleformat{\subsection}
  {\large\bfseries\color{darknavy}}{\thesubsection}{1em}{}
\titleformat{\subsubsection}
  {\normalsize\bfseries\color{darknavy}}{\thesubsubsection}{1em}{}

% ---------- Code listing style ----------
\lstset{
    basicstyle=\ttfamily\small,
    backgroundcolor=\color{lightslate},
    frame=single,
    rulecolor=\color{bordergrey},
    breaklines=true,
    captionpos=b,
}

% ---------- Tip / Note boxes ----------
\newtcolorbox{tipbox}[1][]{
    colback=blue!5, colframe=primaryblue,
    fonttitle=\bfseries, title={#1},
    boxrule=0.5pt, arc=4pt,
}

\newtcolorbox{warnbox}[1][]{
    colback=red!5, colframe=errorred,
    fonttitle=\bfseries, title={#1},
    boxrule=0.5pt, arc=4pt,
}

\newtcolorbox{notebox}[1][]{
    colback=yellow!5, colframe=warningamber,
    fonttitle=\bfseries, title={#1},
    boxrule=0.5pt, arc=4pt,
}

% =============================================================================
\begin{document}

% -------- Title page --------
\begin{titlepage}
\centering
\vspace*{3cm}

{\Huge\bfseries\color{darknavy} Bubble Radius Fitting}\\[0.5cm]
{\LARGE\color{primaryblue} User Guide}\\[2cm]
{\Large Version 1.0}\\[1cm]
{\large Zixiang (Zach) Tong}\\[0.3cm]
{\large The University of Texas at Austin}\\[2cm]
{\large \today}

\vfill
{\small
    Built with Python 3.10+ \quad|\quad PyQt6 \quad|\quad OpenCV \quad|\quad NumPy
}
\end{titlepage}

% -------- Table of contents --------
\tableofcontents
\newpage

% =============================================================================
\section{Introduction}
% =============================================================================

Bubble Radius Fitting is a desktop application for extracting bubble radius
versus time data from high-speed camera image sequences. The software detects
bubble boundaries via adaptive thresholding and morphological filtering, then
fits circles using the Taubin algebraic method.

\subsection{Key Features}

\begin{itemize}[leftmargin=1.5em]
    \item Three analysis modes: Pre-tune, Manual, and Automatic
    \item Side-by-side image viewer with zoom, pan, and overlay visualisations
    \item Interactive ROI (Region of Interest) selection
    \item Advanced image filters: Gaussian Blur, CLAHE, Morphological Open/Close
    \item Live-updating $R(t)$ scatter chart
    \item Data export to \texttt{.mat} files (MATLAB-compatible)
    \item Saveable/loadable parameter presets (JSON)
    \item Radius outlier filtering
\end{itemize}

\subsection{Supported Image Formats}

The application supports standard image file formats commonly produced by
high-speed cameras:

\begin{center}
\begin{tabular}{ll}
    \toprule
    \textbf{Format} & \textbf{Extensions} \\
    \midrule
    TIFF (preferred) & \texttt{.tiff}, \texttt{.tif} \\
    PNG              & \texttt{.png} \\
    JPEG             & \texttt{.jpg}, \texttt{.jpeg} \\
    BMP              & \texttt{.bmp} \\
    \bottomrule
\end{tabular}
\end{center}

When multiple formats are present in the same folder, TIFF files are
given priority. Files within each format are sorted using natural ordering
(e.g.\ \texttt{img1}, \texttt{img2}, \texttt{img10} rather than
\texttt{img1}, \texttt{img10}, \texttt{img2}).

Both 8-bit and 16-bit images are supported; bit depth is auto-detected from
the first image in the sequence.


% =============================================================================
\section{Installation}
% =============================================================================

\subsection{From Source}

\begin{enumerate}
    \item Ensure Python 3.10 or later is installed.
    \item Clone or download the repository.
    \item Create a virtual environment and install dependencies:
\end{enumerate}

\begin{lstlisting}[language=bash]
cd bubbletrack
python -m venv .venv

# Windows
.venv\Scripts\activate
# macOS / Linux
source .venv/bin/activate

pip install -e .
\end{lstlisting}

\subsection{Running the Application}

\begin{lstlisting}[language=bash]
# Run via console entry point
bubbletrack

# Or run as a Python module
python -m bubbletrack
\end{lstlisting}

\subsection{Standalone Executable}

A pre-built Windows executable (\texttt{BubbleRadiusFitting.exe}) is available
in the \texttt{dist/} directory. No Python installation is required. Simply
double-click the executable to launch.


% =============================================================================
\section{Interface Overview}
% =============================================================================

The application uses a three-panel layout inspired by modern design principles.

\subsection{Layout Structure}

\begin{lstlisting}[basicstyle=\ttfamily\footnotesize, frame=none, backgroundcolor=\color{white}]
+-------------------------------------------------------------+
|  HEADER BAR                                   Status: Ready  |
+----------------+--+-----------------------------------------+
|  LEFT SIDEBAR  |  |  MAIN CONTENT AREA                      |
|  (collapsible) |<>|  +------------------+------------------+ |
|  - Image Source|  |  | Original Image   | Binary Mask      | |
|  - Mode Tabs   |  |  |                  |                  | |
|  - Parameters  |  |  +------------------+------------------+ |
|  - Post Proc.  |  |                                         |
|                |  |  Frame Scrubber: [<] [======o=] [>] 1/N  |
|                |  |                                         |
|                |  |  R(t) Scatter Chart                     |
+----------------+--+-----------------------------------------+
|  STATUS BAR: Frame 1/N | .tif | 3.2 um/px | ROI | Mode     |
+-------------------------------------------------------------+
\end{lstlisting}

\subsection{Header Bar}

The dark header bar displays the application title and a colour-coded status
indicator:

\begin{center}
\begin{tabular}{lll}
    \toprule
    \textbf{Colour} & \textbf{Meaning} & \textbf{Example} \\
    \midrule
    \textcolor{successgreen}{\faCircle} Green  & Success / Ready  & ``Ready'', ``R = 42.3 px'' \\
    \textcolor{warningamber}{\faCircle} Amber  & In progress / Warning & ``Fitting\ldots'', ``Stopping\ldots'' \\
    \textcolor{errorred}{\faCircle} Red    & Error            & ``Too few edge points'' \\
    \bottomrule
\end{tabular}
\end{center}

\subsection{Left Sidebar}

The sidebar (340\,px wide) contains all controls, organised into four sections:

\begin{enumerate}
    \item \textbf{Image Source} --- folder selection and file information
    \item \textbf{Mode Tabs} --- Pre-tune, Manual, Automatic
    \item \textbf{Mode-specific Controls} --- parameters and action buttons
    \item \textbf{Post Processing} --- export controls and physical unit settings
\end{enumerate}

The sidebar can be collapsed by clicking the toggle button (\texttt{<}/\texttt{>})
between the sidebar and the main content area, providing more space for the
image panels.

\subsection{Image Panels}

Two image panels are displayed side by side:

\begin{itemize}
    \item \textbf{Original Image} (left) --- the greyscale source image with
          overlay graphics: green ROI rectangle, blue fitted circle, and red
          edge detection points.
    \item \textbf{Binary Mask} (right) --- the processed binary image showing
          the detection result.
\end{itemize}

Each panel provides zoom controls in the top-right corner:

\begin{center}
\begin{tabular}{cl}
    \toprule
    \textbf{Control} & \textbf{Function} \\
    \midrule
    \texttt{[+]}  & Zoom in ($1.25\times$) \\
    \texttt{[-]}  & Zoom out ($0.8\times$) \\
    \texttt{[$\circlearrowleft$]} & Reset zoom to fit the view \\
    Scroll wheel  & Zoom in/out ($1.15\times$ per notch, anchored at cursor) \\
    \bottomrule
\end{tabular}
\end{center}

\subsection{Frame Scrubber}

The frame scrubber bar sits between the image panels and the chart. It
contains:

\begin{itemize}
    \item \textbf{Previous/Next buttons} (\texttt{<}/\texttt{>}) --- step one
          frame backward or forward.
    \item \textbf{Slider} --- drag to jump to any frame in the sequence.
    \item \textbf{Frame label} --- displays the current frame number and total
          (e.g.\ ``Frame: 42/500'').
\end{itemize}

\subsection{Radius Chart}

A Matplotlib scatter plot showing fitted radius (pixels) versus frame number.
Data points appear as red ``+'' markers. The chart updates in real time during
fitting operations.

\subsection{Status Bar}

The bottom status bar displays contextual information:

\begin{center}
\begin{tabular}{ll}
    \toprule
    \textbf{Field} & \textbf{Example} \\
    \midrule
    Frame counter  & Frame: 42/500 \\
    Image format   & Format: .tif \\
    Physical scale & Scale: 3.2 \textmu m/px \\
    ROI bounds     & ROI: [1--512, 1--512] \\
    Current mode   & Mode: Pre-tune \\
    \bottomrule
\end{tabular}
\end{center}


% =============================================================================
\section{Quick Start}
% =============================================================================

\begin{enumerate}
    \item \textbf{Load images:} Click \textbf{Browse} in the Image Source
          section and select a folder containing image files.

    \item \textbf{Navigate frames:} Use the frame scrubber slider or the
          \texttt{<}/\texttt{>} buttons to browse the image sequence.

    \item \textbf{Select ROI:} In the \textbf{Pre-tune} tab, click
          \textbf{Select ROI}, then drag a rectangle on the original image
          to define the region of interest.

    \item \textbf{Tune parameters:} Adjust the \textbf{Threshold} and
          \textbf{Removing Factor} sliders until the bubble boundary appears
          cleanly in the binary mask. Enable advanced filters if needed.

    \item \textbf{Fit a single frame:} Click \textbf{Fit Current Frame}. The
          fitted circle (blue) and edge points (red) appear on the original
          image.

    \item \textbf{Batch process:} Switch to the \textbf{Automatic} tab, set
          the frame range, and click \textbf{Fit All}. Progress is shown in
          real time.

    \item \textbf{Export results:} In the \textbf{Post Processing} section,
          set the physical parameters (FPS, scale) and click
          \textbf{Export Pixel Data} or \textbf{Export Physical Data} to save
          results as \texttt{.mat} files.
\end{enumerate}


% =============================================================================
\section{Analysis Modes}
% =============================================================================

% ---------- Pre-tune ----------
\subsection{Pre-tune Mode}

Pre-tune mode is designed for parameter optimisation on a single frame. All
parameter changes update the binary mask preview in real time (debounced at
50\,ms), allowing rapid visual feedback.

\subsubsection{Frame Selector}

A spin box allows direct navigation to any frame in the sequence (1-indexed
display). This is independent of the main frame scrubber.

\subsubsection{Region of Interest (ROI)}

\begin{enumerate}
    \item Click \textbf{Select ROI}. The image border turns amber to indicate
          selection mode.
    \item Drag a rectangle on the original image panel.
    \item The ROI coordinates are displayed in the sidebar and the status bar.
\end{enumerate}

By default, the ROI covers the entire image. Only pixels within the ROI are
used for bubble detection and circle fitting. The ROI is shown as a green
rectangle on the original image.

\begin{tipbox}[Tip]
    Choose an ROI that tightly encloses the bubble. A smaller ROI improves
    detection accuracy and processing speed.
\end{tipbox}

\subsubsection{Threshold}

\begin{center}
\begin{tabular}{llll}
    \toprule
    \textbf{Range} & \textbf{Default} & \textbf{Unit} & \textbf{Effect} \\
    \midrule
    0--100 & 50 & \% & Adaptive binarisation sensitivity \\
    \bottomrule
\end{tabular}
\end{center}

The threshold slider controls the sensitivity of adaptive mean thresholding.
Higher values produce more white (foreground) pixels; lower values produce
a stricter foreground detection. The internal sensitivity value is computed as:
\[
    \text{sensitivity} = \frac{\text{slider value}}{100}
\]

\subsubsection{Removing Factor}

\begin{center}
\begin{tabular}{llll}
    \toprule
    \textbf{Range} & \textbf{Default} & \textbf{Unit} & \textbf{Effect} \\
    \midrule
    0--100 & 90 & \% & Minimum connected-component area (log scale) \\
    \bottomrule
\end{tabular}
\end{center}

This slider removes small white speckles (noise) from the binary image.
The mapping is logarithmic: a slider value of 0 removes components smaller
than 10\,px$^2$; a slider value of 100 removes components up to 50\%\ of
the ROI area.

\subsubsection{Edge Crossing Flags}

Four checkboxes (Top, Right, Down, Left) indicate which edges of the ROI the
bubble extends beyond. When a flag is checked, the algorithm expands the
binary image in the \emph{opposite} direction to correctly detect partial
bubbles.

\begin{notebox}[Note]
    For example, if the bubble extends beyond the top and right edges of the
    ROI, check \textbf{Top} and \textbf{Right}. The algorithm will expand
    the image downward and leftward to compensate.
\end{notebox}

\subsubsection{Advanced Filters}

The Advanced Filters section (initially collapsed) provides four optional
image filters. Enable each filter via its checkbox.

\begin{center}
\begin{tabular}{lllp{6cm}}
    \toprule
    \textbf{Filter} & \textbf{Range} & \textbf{Default} & \textbf{Description} \\
    \midrule
    Gaussian Blur & 1--20 & 3 & Applied \emph{before} binarisation.
        Smooths the greyscale image to suppress small noise such as tracer
        particles. \\[4pt]
    CLAHE & 1--40 & 2 & Applied \emph{before} binarisation. Contrast
        Limited Adaptive Histogram Equalisation improves local contrast
        when lighting is uneven. \\[4pt]
    Morph Close & 1--30 & 5 & Applied \emph{after} binarisation. Fills
        small holes and gaps inside the bubble boundary. \\[4pt]
    Morph Open & 1--20 & 3 & Applied \emph{after} binarisation. Removes
        thin protrusions and spurs on the bubble edge. \\
    \bottomrule
\end{tabular}
\end{center}

\subsubsection{Parameter Presets}

Click \textbf{Save Preset} to store the current parameter configuration
(threshold, removing factor, edge flags, and all filter settings) to a JSON
file. Click \textbf{Load Preset} to restore a previously saved configuration.

\subsubsection{Fit Current Frame}

Click \textbf{Fit Current Frame} to run the full detection and circle-fitting
pipeline on the current frame. On success:

\begin{itemize}
    \item A \textcolor{primaryblue}{blue circle} is drawn on the original
          image showing the fitted bubble boundary.
    \item \textcolor{errorred}{Red dots} mark the detected edge points.
    \item The header displays the fitted radius in pixels.
    \item The $R(t)$ chart is updated.
\end{itemize}

% ---------- Manual ----------
\subsection{Manual Mode}

Manual mode allows you to select edge points by clicking directly on the
bubble boundary.

\begin{enumerate}
    \item Click \textbf{Select Points}. The image border turns blue and the
          cursor changes to a crosshair.
    \item Click on the bubble edge to place points. Each click adds a red
          marker on the image. The point counter updates in real time.
    \item A minimum of \textbf{3 points} is required. The status indicator
          turns green when enough points are collected.
    \item Click \textbf{Done} to fit a circle through the selected points
          using the Taubin algorithm.
    \item Click \textbf{Clear} to discard all points and start over.
\end{enumerate}

\begin{tipbox}[Tip]
    For best results, distribute the selected points evenly around the
    bubble perimeter. Using 5--10 points typically gives reliable fits.
\end{tipbox}

% ---------- Automatic ----------
\subsection{Automatic Mode}

Automatic mode batch-processes a range of frames using the parameters
configured in the Pre-tune tab.

\subsubsection{Setting the Frame Range}

Use the \textbf{From} and \textbf{To} spin boxes to define the processing
range (1-indexed, inclusive). By default, the range covers all frames.

\subsubsection{Running Batch Processing}

\begin{enumerate}
    \item Click \textbf{Fit All} to start processing. A background worker
          thread processes each frame sequentially.
    \item The progress bar and percentage label update in real time.
    \item The original image, binary mask, and $R(t)$ chart are updated
          periodically during processing (${\sim}150$\,ms intervals).
    \item The fitted circle and edge points are displayed on the original
          image for each processed frame.
\end{enumerate}

\subsubsection{Stopping}

Click \textbf{Stop} to request a graceful stop. The worker finishes the
current frame and then halts. All results computed so far are retained.

\subsubsection{Clearing Results}

Click \textbf{Clear} to discard all fitting results and reset the $R(t)$
chart. This does not affect the loaded images or parameter settings.

\begin{warnbox}[Warning]
    Clicking \textbf{Clear} permanently removes all fitting results from
    memory. Export your data before clearing if you need to keep the results.
\end{warnbox}


% =============================================================================
\section{Post Processing and Export}
% =============================================================================

The Post Processing section (in the lower portion of the left sidebar)
provides physical unit conversion parameters and two export options.

\subsection{Physical Parameters}

\begin{center}
\begin{tabular}{llllp{5.5cm}}
    \toprule
    \textbf{Parameter} & \textbf{Range} & \textbf{Default} & \textbf{Unit} & \textbf{Description} \\
    \midrule
    FPS & 0.001--99\,999 & 10 & Hz, kHz, or MHz & Frame rate of the camera.
        Select the unit from the dropdown (Hz, k, M). \\[4pt]
    Scale & 0.001--9\,999 & 3.2 & \textmu m/px & Spatial calibration factor. \\[4pt]
    Rmax Fit Length & 3--999 (odd) & 11 & frames & Number of frames used for
        quadratic $R_{\max}$ interpolation (must be odd). \\
    \bottomrule
\end{tabular}
\end{center}

\subsection{Export Pixel Data}

Click \textbf{Export Pixel Data} to save raw fitting results. A file dialog
opens with the default filename \texttt{radius\_pixel.mat} in the image
folder. The \texttt{.mat} file contains:

\begin{center}
\begin{tabular}{llp{7cm}}
    \toprule
    \textbf{Variable} & \textbf{Size} & \textbf{Description} \\
    \midrule
    \texttt{Radius} & $(1 \times N)$ & Fitted radius for each frame in pixels.
        Unfitted frames contain $-1$. \\
    \texttt{CircleFitPar} & $(N \times 2)$ & Circle centre coordinates
        $[\text{row},\; \text{col}]$ for each frame. \\
    \texttt{CircleXY} & $\{1 \times N\}$ & Cell array of edge-point
        coordinate matrices $(M_i \times 2)$ for each frame. \\
    \bottomrule
\end{tabular}
\end{center}

\subsection{Export Physical Data}

Click \textbf{Export Physical Data} to save time-resolved radius data in
physical units. A file dialog opens with the default filename
\texttt{radius\_time\_physical.mat}. The export process:

\begin{enumerate}
    \item Converts pixel radii to physical units using the scale factor.
    \item Converts frame indices to time using the frame rate.
    \item Identifies the frame with maximum radius.
    \item Fits a quadratic curve $R(t) = at^2 + bt + c$ around the peak
          using the specified fit window.
    \item If the parabola is concave ($a < 0$), computes the sub-frame
          peak: $t_{\text{peak}} = -b/(2a)$, $R_{\text{peak}} = (4ac - b^2)/(4a)$.
    \item Shifts the time axis so $t = 0$ at the peak radius location.
\end{enumerate}

The \texttt{.mat} file contains:

\begin{center}
\begin{tabular}{llp{7cm}}
    \toprule
    \textbf{Variable} & \textbf{Size} & \textbf{Description} \\
    \midrule
    \texttt{t} & $(N{+}1 \times 1)$ & Time array in seconds, centred at
        $R_{\max}$. \\
    \texttt{R} & $(N{+}1 \times 1)$ & Radius array in \textmu m. \\
    \texttt{RmaxAll} & scalar & Peak radius value in \textmu m. \\
    \texttt{RmaxTimeLoc} & scalar & Time of peak radius in seconds. \\
    \bottomrule
\end{tabular}
\end{center}

\begin{notebox}[Note]
    If the $R_{\max}$ fit window extends beyond the data boundaries, the
    export will fail with an error message. Reduce the \textbf{Rmax Fit
    Length} parameter or ensure the peak radius is sufficiently far from
    the edges of the frame range.
\end{notebox}


% =============================================================================
\section{Algorithm Overview}
% =============================================================================

This section provides a technical description of the image processing and
circle-fitting pipeline.

\subsection{Processing Pipeline}

The detection pipeline consists of the following steps, executed in order:

\begin{enumerate}
    \item \textbf{Load and normalise} --- Read the raw image and normalise
          pixel values to the $[0, 1]$ range based on the detected bit depth.

    \item \textbf{Gaussian Blur} (optional) --- Apply a Gaussian smoothing
          kernel with the specified $\sigma$ to suppress high-frequency noise.
          Applied \emph{before} binarisation.

    \item \textbf{CLAHE} (optional) --- Apply Contrast Limited Adaptive
          Histogram Equalisation with the specified clip limit to improve
          local contrast. Applied \emph{before} binarisation.

    \item \textbf{Adaptive thresholding} --- Convert the greyscale image to
          binary using local mean thresholding (OpenCV
          \texttt{adaptiveThreshold}). The block size is computed
          automatically as:
          \[
              B = 2 \left\lceil \frac{\min(H, W)}{16} \right\rceil + 1
              \qquad (B \geq 3)
          \]
          The threshold offset is derived from the sensitivity parameter:
          \[
              C = (\text{sensitivity} - 0.5) \times 255
          \]

    \item \textbf{ROI extraction} --- Crop the binary image to the
          user-defined region of interest.

    \item \textbf{Boundary expansion} --- If the bubble extends beyond the
          ROI, the image is doubled along the non-crossing edge direction
          and filled with white pixels to complete partial bubbles.

    \item \textbf{Small object removal} --- Connected components with area
          below the removing factor threshold are removed.

    \item \textbf{Morphological opening} (optional) --- Erosion followed by
          dilation with a disk structuring element to remove thin protrusions.

    \item \textbf{Morphological closing} (optional) --- Dilation followed by
          erosion with a disk structuring element, then hole filling, to close
          gaps in the bubble boundary.

    \item \textbf{Inversion} --- The image is inverted so that the bubble
          interior (dark in the original) becomes the foreground.

    \item \textbf{Blob filtering} --- Connected components are labelled and
          filtered by shape:
          \begin{itemize}
              \item Reject if major axis $> 2.2 \times$ short image side
              \item Reject if eccentricity $> 1.6$ (major/minor axis ratio)
              \item Select the largest remaining blob
          \end{itemize}

    \item \textbf{Edge detection} --- Extract the boundary using a
          morphological gradient: dilate the blob mask with a $3 \times 3$
          kernel, then XOR with the original mask.

    \item \textbf{Crop back} --- If the image was expanded in step~6, crop
          back to the original ROI dimensions.
\end{enumerate}

\subsection{Circle Fitting: Taubin Method}

The application uses the Taubin algebraic method for circle fitting, which
provides a robust least-squares estimate that is invariant to rotation and
translation.

Given $N \geq 3$ boundary points $\{(x_i, y_i)\}$, the algorithm:

\begin{enumerate}
    \item Computes the centroid $(\bar{x}, \bar{y})$ and centres the data.
    \item Builds a covariance-based system from the centred coordinates.
    \item Solves a characteristic polynomial via Newton's method (up to 20
          iterations, tolerance $10^{-12}$).
    \item Extracts the circle centre $(x_c, y_c)$ and radius $r$ from the
          solution.
\end{enumerate}

\begin{tipbox}[Reference]
    G.\ Taubin, ``Estimation Of Planar Curves, Surfaces And Nonplanar Space
    Curves Defined By Implicit Equations, With Applications To Edge And Range
    Image Segmentation,'' \emph{IEEE Trans.\ PAMI}, Vol.~13, pp.~1115--1138,
    1991.
\end{tipbox}

\subsection{Radius Outlier Rejection}

After circle fitting, the fitted radius is compared against the maximum image
dimension (the longer side of the image). If the radius exceeds this value or
is \texttt{NaN}, the result is rejected and the frame is marked as unfitted.

\subsection{$R_{\max}$ Quadratic Interpolation}

During physical data export, the software identifies the maximum radius and
fits a parabola to the surrounding data points:

\begin{equation}
    R(t) = at^2 + bt + c
\end{equation}

If the fit yields a concave parabola ($a < 0$), the sub-frame peak is
computed analytically:
\begin{align}
    t_{\text{peak}} &= -\frac{b}{2a} \\
    R_{\text{peak}} &= \frac{4ac - b^2}{4a}
\end{align}

This provides sub-frame temporal precision for determining the maximum bubble
radius.


% =============================================================================
\section{Parameter Reference}
% =============================================================================

\subsection{Image Processing Parameters}

\begin{center}
\begin{tabular}{llllp{4.5cm}}
    \toprule
    \textbf{Parameter} & \textbf{Range} & \textbf{Default} & \textbf{Unit}
        & \textbf{Effect} \\
    \midrule
    Threshold & 0--100 & 50 & \% & Adaptive binarisation sensitivity.
        Higher = more foreground pixels. \\[3pt]
    Removing Factor & 0--100 & 90 & \% & Minimum component area
        (logarithmic scale). Higher = removes larger noise. \\[3pt]
    Gaussian $\sigma$ & 1--20 & 3 & px & Pre-binarisation smoothing
        strength. \\[3pt]
    CLAHE Clip & 1--40 & 2 & --- & Contrast enhancement strength. \\[3pt]
    Morph Close & 1--30 & 5 & px & Post-binarisation gap-filling
        radius. \\[3pt]
    Morph Open & 1--20 & 3 & px & Post-binarisation spur-removal
        radius. \\[3pt]
    Edge Flags & 4 bools & all off & --- & Bubble crosses
        [Top, Right, Down, Left] edges. \\
    \bottomrule
\end{tabular}
\end{center}

\subsection{Physical Conversion Parameters}

\begin{center}
\begin{tabular}{lllll}
    \toprule
    \textbf{Parameter} & \textbf{Range} & \textbf{Default} & \textbf{Unit}
        & \textbf{Effect} \\
    \midrule
    FPS & 0.001--99\,999 & 10 & Hz/kHz/MHz & Camera frame rate \\
    Scale & 0.001--9\,999 & 3.2 & \textmu m/px & Pixel calibration \\
    Rmax Fit Length & 3--999 & 11 & frames & Quadratic fit window (odd) \\
    \bottomrule
\end{tabular}
\end{center}

\subsection{Internal Constraints}

\begin{center}
\begin{tabular}{lll}
    \toprule
    \textbf{Constraint} & \textbf{Value} & \textbf{Description} \\
    \midrule
    Min.\ edge points & 3 & Required for circle fitting \\
    Max.\ radius & $\max(H, W)$ & Image long-side dimension \\
    Min.\ hole area & 40\,px$^2$ & Holes smaller than this are filled \\
    Max.\ axis ratio & 2.2 & Blob major axis / short image side \\
    Max.\ eccentricity & 1.6 & Major / minor axis ratio \\
    \bottomrule
\end{tabular}
\end{center}


% =============================================================================
\section{Typical Workflows}
% =============================================================================

\subsection{Workflow 1: Single-Frame Analysis}

This workflow is useful for quickly inspecting individual frames or verifying
parameter settings.

\begin{enumerate}
    \item Load the image folder via \textbf{Browse}.
    \item Navigate to the frame of interest using the frame scrubber.
    \item In the \textbf{Pre-tune} tab:
    \begin{enumerate}[label=\alph*.]
        \item Select an ROI around the bubble.
        \item Adjust \textbf{Threshold} until the bubble interior is clearly
              segmented.
        \item Increase \textbf{Removing Factor} to eliminate noise spots.
        \item Enable \textbf{Gaussian Blur} if small particles (e.g.\ tracer
              beads) are present.
        \item Enable \textbf{CLAHE} if lighting is uneven.
        \item Check edge flags if the bubble is partially outside the ROI.
    \end{enumerate}
    \item Click \textbf{Fit Current Frame} and inspect the result.
    \item Optionally save the parameter preset for later use.
\end{enumerate}

\subsection{Workflow 2: Batch Processing}

This is the primary workflow for extracting $R(t)$ data from a complete image
sequence.

\begin{enumerate}
    \item Complete the parameter tuning (Workflow~1) on a representative frame.
    \item Switch to the \textbf{Automatic} tab.
    \item Set the \textbf{From} and \textbf{To} frame numbers.
    \item Click \textbf{Fit All}. Monitor progress via the progress bar and
          the $R(t)$ chart.
    \item Once complete, review the chart for outliers or missing data points.
    \item In \textbf{Post Processing}:
    \begin{enumerate}[label=\alph*.]
        \item Set the camera \textbf{FPS} and spatial \textbf{Scale}.
        \item Click \textbf{Export Pixel Data} for raw results.
        \item Click \textbf{Export Physical Data} for time-resolved
              physical-unit results with $R_{\max}$ interpolation.
    \end{enumerate}
\end{enumerate}

\subsection{Workflow 3: Manual Point Selection}

Use this workflow when the automatic detection fails to identify the bubble
boundary correctly (e.g.\ due to poor contrast or complex shapes).

\begin{enumerate}
    \item Navigate to the target frame.
    \item Switch to the \textbf{Manual} tab.
    \item Click \textbf{Select Points} and click 5--10 points around the
          bubble edge.
    \item Click \textbf{Done} to fit the circle.
    \item Repeat for additional frames if needed.
    \item Export results as in Workflow~2.
\end{enumerate}


% =============================================================================
\section{Troubleshooting}
% =============================================================================

\begin{center}
\renewcommand{\arraystretch}{1.4}
\begin{tabular}{p{4.5cm}p{4cm}p{5cm}}
    \toprule
    \textbf{Problem} & \textbf{Possible Cause} & \textbf{Solution} \\
    \midrule

    ``No images found'' after selecting a folder &
    Folder contains no supported image files &
    Ensure the folder contains \texttt{.tiff}, \texttt{.tif}, \texttt{.png},
    \texttt{.jpg}, or \texttt{.bmp} files. \\

    Binary mask is mostly white or mostly black &
    Threshold too high or too low &
    Adjust the Threshold slider until the bubble boundary is clearly
    visible. \\

    Small noise spots remain in the binary mask &
    Removing Factor too low &
    Increase the Removing Factor slider. Enable Gaussian Blur to suppress
    fine noise. \\

    ``Too few edge points'' after fitting &
    Bubble not detected or ROI too small &
    Enlarge the ROI, adjust the Threshold, or check the edge flags. \\

    ``Radius outlier, skipped'' &
    Fitted circle is unreasonably large &
    The detection picked up the wrong region. Adjust Threshold, Removing
    Factor, or use Morph Open/Close to clean the boundary. \\

    Bubble boundary has gaps &
    Weak contrast at some edges &
    Enable \textbf{CLAHE} to improve local contrast, or enable
    \textbf{Morph Close} to fill gaps. \\

    Bubble boundary has spurs or protrusions &
    Noise connected to the boundary &
    Enable \textbf{Morph Open} to remove thin protrusions. \\

    Stop button does not respond immediately &
    Worker finishes current frame before stopping &
    This is expected. The stop is graceful---wait a moment for the
    current frame to complete. \\

    Export fails with ``Rmax\_Fit\_Length exceeds data range'' &
    Fit window is wider than available data around the peak &
    Reduce the \textbf{Rmax Fit Length} parameter, or process more
    frames around the peak radius. \\

    Application window is too small &
    Display scaling or resolution issue &
    The window can be freely resized. Collapse the left sidebar for
    more image space. \\

    \bottomrule
\end{tabular}
\end{center}


% =============================================================================
\section{Keyboard and Mouse Reference}
% =============================================================================

\begin{center}
\renewcommand{\arraystretch}{1.3}
\begin{tabular}{llll}
    \toprule
    \textbf{Mode} & \textbf{Input} & \textbf{Action} & \textbf{Result} \\
    \midrule
    Any & Scroll wheel (on image) & Zoom in/out & $1.15\times$ per notch \\
    Any & Click \texttt{[+]}/\texttt{[-]}/\texttt{[$\circlearrowleft$]}
        & Zoom/Reset & Button-based zoom control \\
    Normal & Drag on image & Pan & Scroll hand drag \\
    Point  & Left-click on image & Place point & Red dot added \\
    ROI    & Left-drag on image & Draw rectangle & ROI defined \\
    Any & Click \texttt{<}/\texttt{>} & Step frame & $\pm 1$ frame \\
    Any & Drag frame slider & Jump to frame & Instant navigation \\
    \bottomrule
\end{tabular}
\end{center}


% =============================================================================
\section{Export File Format Details}
% =============================================================================

All exported files use the MATLAB \texttt{.mat} (v5) format and can be loaded
in MATLAB using the \texttt{load} function or in Python using
\texttt{scipy.io.loadmat}.

\subsection{Pixel Data (\texttt{radius\_pixel.mat})}

\begin{lstlisting}[language=Matlab]
data = load('radius_pixel.mat');

% Radius array (1 x N), in pixels. -1 = unfitted frame.
R = data.Radius;

% Circle centres (N x 2), [row, col] for each frame.
centres = data.CircleFitPar;

% Edge points: cell array {1 x N}.
% Each cell contains an (M x 2) matrix of [row, col] coordinates.
edges = data.CircleXY;
\end{lstlisting}

\subsection{Physical Data (\texttt{radius\_time\_physical.mat})}

\begin{lstlisting}[language=Matlab]
data = load('radius_time_physical.mat');

% Time array (N+1 x 1), in seconds. Centred at Rmax.
t = data.t;

% Radius array (N+1 x 1), in micrometres.
R = data.R;

% Maximum radius (scalar), in micrometres.
Rmax = data.RmaxAll;

% Time of maximum radius (scalar), in seconds.
t_Rmax = data.RmaxTimeLoc;

% Plot R(t)
figure;
plot(t * 1e6, R, 'r.');  % convert to microseconds
xlabel('Time (\mus)');
ylabel('Radius (\mum)');
title(sprintf('R_{max} = %.1f \\mum', Rmax));
\end{lstlisting}


% =============================================================================
\section*{Acknowledgements}
% =============================================================================
\addcontentsline{toc}{section}{Acknowledgements}

This software was developed at The University of Texas at Austin. The circle
fitting algorithm is based on the Taubin method as described in:

\begin{quote}
    G.\ Taubin, ``Estimation Of Planar Curves, Surfaces And Nonplanar Space
    Curves Defined By Implicit Equations, With Applications To Edge And Range
    Image Segmentation,'' \emph{IEEE Trans.\ PAMI}, Vol.~13, pp.~1115--1138,
    1991.
\end{quote}

% =============================================================================
\end{document}
